\section{Introduction}
\label{sec:intro}

Permissioned Byzantine Fault Tolerant (BFT) consensus underpins e-governance
systems --- voting, registries, digital identity --- and forces a concrete
procurement question: \emph{how many validators should the network have?}

The validator count~$\n$ pulls in opposite directions. Throughput argues for
small~$\n$: every validator executes every transaction, and PBFT-family
protocols exchange $O(\n^{2})$ messages per block. Detection reliability argues
for large~$\n$. Combining the two for sizing is the subject of this paper.

Published single-factor fits $\TPS(\n)\propto\n^{\alpha}$ disagree in magnitude
and sometimes in sign. Our earlier studies, extrapolated to sixty validators,
differ by roughly an order of magnitude. This disagreement has often been
interpreted as protocol-specific behaviour. We show that it can arise solely
from measurement design: concurrency is rarely held fixed while~$\n$ varies, so
the fitted scaling exponent mixes load-generator growth with consensus
overhead. Without reporting the concurrency path, published exponents cannot be
read as protocol properties. Packing many validators onto one host adds a second
confound --- CPU contention --- that resource-normalized measurement isolates.

\paragraph{Contributions.}
On the methodological front, we introduce a resource-normalized measurement
(\RNM{}) protocol that fixes CPU quota~$\q$ per validator and reports
$\TPSn=\TPS/\q$, isolating intrinsic consensus cost from host contention ---
analogous to a calibrated wind tunnel. Analytically, along a concurrency path
of slope~$\lambda$ the single-factor OLS fit satisfies
$\ahat_{m}\xrightarrow{\Prob}\gamma\lambda-\beta$
(Proposition~\ref{prop:identifiability}); the value of the result is the
engineering interpretation that reconciles historical BFT benchmarks.
From a planning perspective, we formulate chance-constrained sizing whose
feasible set is $[\nmin,\nmax]$ (Proposition~\ref{prop:window}). Empirically, we
validate prediction intervals by coverage and run a public Besu/QBFT CI
campaign (Sections~\ref{sec:exp}--\ref{sec:case}). Scope: throughput and
detection; absolute TPS environment-specific; cross-client replication out of
scope.
